\subsubsection{Alice's Operations and Mathematical Derivation}

The teleportation protocol applies two gates to the first two qubits of the three-qubit state $\ket{\Psi_1}$ (Alice's qubit and her half of the Bell pair). The first gate is a CNOT gate, which uses Alice's qubit as the control and her half of the Bell pair as the target. The second gate is a Hadamard gate applied to Alice's qubit. Let's explain the effect of these gates on the state $\ket{\Psi_1}$.

\highspace
\begin{flushleft}
    \textcolor{Green3}{\faIcon{tools} \textbf{Apply the CNOT}}
\end{flushleft}
The \textbf{first operation is a CNOT gate applied to the first two qubits}. The \hl{control is the first qubit $q_{\psi}$} and the \hl{target is the second qubit $q_{A}$.} So the second qubit flips only when the first qubit is $1$. Therefore, \hl{only the terms $\beta \ket{100}$ and $\beta \ket{111}$ will be affected} by the CNOT gate. The first term $\beta \ket{100}$ will become $\beta \ket{110}$, while the second term $\beta \ket{111}$ will become $\beta \ket{101}$. The other terms remain unchanged.
\begin{equation}
    \ket{\Psi_{\text{CX}}} = \frac{1}{\sqrt{2}} \left(
        \alpha \ket{000} +
        \alpha \ket{011} +
        \beta \ket{110} +
        \beta \ket{101}
    \right)
\end{equation}
\begin{center}
    \begin{quantikz}
        \lstick{$q_{\psi}$} & \ctrl{1}  &  \\
        \lstick{$q_{A}$}    & \targ{}   &  \\
        \lstick{$q_{B}$}    &           & 
    \end{quantikz}
\end{center}

\noindent
Notice that \hl{Bob's qubit, the third one, is never touched by Alice's operations, so it remains unchanged.} The state of the system after the CNOT gate is $\ket{\Psi_{\text{CX}}}$.

\highspace
\begin{flushleft}
    \textcolor{Green3}{\faIcon{tools} \textbf{Apply the Hadamard to the first qubit}}
\end{flushleft}
Now \textbf{Alice applies a Hadamard gate to her qubit} $q_{\psi}$. Recall that the Hadamard gate transforms the basis states as follows:
\begin{equation*}
    H \ket{0} = \frac{1}{\sqrt{2}} (\ket{0} + \ket{1}), \quad
    H \ket{1} = \frac{1}{\sqrt{2}} (\ket{0} - \ket{1})
\end{equation*}
For each term in $\ket{\Psi_{\text{CX}}}$, we apply the Hadamard transformation to the first qubit. So, for convenience, we can rewrite the state $\ket{\Psi_{\text{CX}}}$ as:
\begin{equation*}
    \ket{\Psi_{\text{CX}}} = \frac{1}{\sqrt{2}} \left(
        \alpha \ket{0} \otimes (\ket{00} + \ket{11}) +
        \beta \ket{1} \otimes (\ket{10} + \ket{01})
    \right)
\end{equation*}
Now we apply the Hadamard gate to the first qubit:
\begin{align*}
    H \ket{\Psi_{\text{CX}}} &= \frac{1}{\sqrt{2}} \left(
        \alpha H\ket{0} \otimes (\ket{00} + \ket{11}) +
        \beta H\ket{1} \otimes (\ket{10} + \ket{01})
    \right) \\
    &= \frac{1}{\sqrt{2}} \left(
        \alpha \frac{1}{\sqrt{2}} (\ket{0} + \ket{1}) \otimes (\ket{00} + \ket{11}) + \right. \\
        & \hspace{1.2cm} \left. \beta \frac{1}{\sqrt{2}} (\ket{0} - \ket{1}) \otimes (\ket{10} + \ket{01})
    \right) \\
    &= \frac{1}{2} \left(
        \alpha (\ket{0} + \ket{1}) \otimes (\ket{00} + \ket{11}) +
        \beta (\ket{0} - \ket{1}) \otimes (\ket{10} + \ket{01})
    \right)
\end{align*}
Expanding this expression, we get:
\begin{align*}
    \ket{\Psi_{\text{H}}} &= \frac{1}{2} \left(
        \alpha \ket{0} \ket{00} + \alpha \ket{0} \ket{11} + \alpha \ket{1} \ket{00} + \alpha \ket{1} \ket{11} + \right. \\
        & \hspace{0.83cm} \left. \beta \ket{0} \ket{10} + \beta \ket{0} \ket{01} - \beta \ket{1} \ket{10} - \beta \ket{1} \ket{01}
    \right) \\
    &= \frac{1}{2} \left(
        \alpha \ket{000} + \alpha \ket{011} + \alpha \ket{100} + \alpha \ket{111} + \right. \\
        & \hspace{0.83cm} \left. \beta \ket{010} + \beta \ket{001} - \beta \ket{110} - \beta \ket{101}
    \right)
\end{align*}
\begin{center}
    \begin{quantikz}
        \lstick{$q_{\psi}$} & \ctrl{1}  & \gate{H} &  \\
        \lstick{$q_{A}$}    & \targ{}   &          &  \\
        \lstick{$q_{B}$}    &           &          & 
    \end{quantikz}
\end{center}
Even if it is correct, it \textbf{does not reveal teleportation}. To see teleportation, we can \hl{rearrange the terms by grouping them according to the owner of the first two qubits (Alice's) and the last qubit (Bob's)}:
\begin{itemize}
    \item Alice's bits $00$:
    \begin{equation*}
        \alpha\ket{000} + \beta\ket{001}
    \end{equation*}
    \begin{equation*}
        \ket{00} \otimes \underbrace{\left(\alpha\ket{0} + \beta\ket{1}\right)}_{\ket{\psi}}
    \end{equation*}


    \item Alice's bits $01$:
    \begin{equation*}
        \alpha\ket{011} + \beta\ket{010}
    \end{equation*}
    \begin{equation*}
        \ket{01} \otimes \left(\alpha\ket{1} + \beta\ket{0}\right) \xrightarrow{\text{convention}} \ket{01} \otimes \left(\beta\ket{0} + \alpha\ket{1}\right)
    \end{equation*}


    \item Alice's bits $10$:
    \begin{equation*}
        \alpha\ket{100} - \beta\ket{101}
    \end{equation*}
    \begin{equation*}
        \ket{10} \otimes \left(\alpha\ket{0} - \beta\ket{1}\right)
    \end{equation*}


    \item Alice's bits $11$:
    \begin{equation*}
        \alpha\ket{111} - \beta\ket{110}
    \end{equation*}
    \begin{equation*}
        \ket{11} \otimes \left(\alpha\ket{1} - \beta\ket{0}\right) \xrightarrow{\text{convention}} \ket{11} \otimes \left(-\beta\ket{0} + \alpha\ket{1}\right)
    \end{equation*}
\end{itemize}
Putting all together, we can write the state after Alice's operations as:
\begin{equation}
    \begin{aligned}
        \ket{\Psi_{\text{H}}} = \ket{\Psi_{2}} &= \frac{1}{2} \left[ \,
            \ket{00} \otimes \left(\alpha\ket{0} + \beta\ket{1}\right) + \right. \\
            & \hspace{0.88cm} \ket{01} \otimes \left(\beta\ket{0} + \alpha\ket{1}\right) + \\
            & \hspace{0.88cm} \ket{10} \otimes \left(\alpha\ket{0} - \beta\ket{1}\right) + \\
            & \hspace{0.84cm} \left. \ket{11} \otimes \left(-\beta\ket{0} + \alpha\ket{1}\right)
        \, \right]
    \end{aligned}
\end{equation}

\newpage

\noindent
\textcolor{Green3}{\faIcon{question-circle} \textbf{What does this equation mean?}} At first sight, it may not be clear what this equation means. However, it is crucial to understand that the first two qubits belong to Alice, while the last qubit belongs to Bob (state in parentheses). The equation shows that the \hl{state of Bob's qubit depends on the measurement outcome of Alice's qubits.}

\highspace
In other words, the equation says:
\begin{itemize}
    \item If Alice later measures $00$, Bob has:
    \begin{equation*}
        \ket{\psi} = \alpha\ket{0} + \beta\ket{1}
    \end{equation*}

    \item If Alice later measures $01$, Bob has:
    \begin{equation*}
        \beta\ket{0} + \alpha\ket{1}
    \end{equation*}

    \item If Alice later measures $10$, Bob has:
    \begin{equation*}
        \alpha\ket{0} - \beta\ket{1}
    \end{equation*}

    \item If Alice later measures $11$, Bob has:
    \begin{equation*}
        -\beta\ket{0} + \alpha\ket{1}
    \end{equation*}
\end{itemize}
So Bob does not always have exactly the state $\ket{\psi}$, but he always has a state related to $\ket{\psi}$ by a simple Pauli gate.

\highspace
In summary, \hl{Alice's two operations have transformed the total system so that Bob's qubit is now 1 of 4 known variants of the original unknown state.} The next section explains how Alice's two classical bits tell Bob which correction to apply.